\section{Problem Formulation and Method}

\subsection{Forecasting Task}

Let
\begin{equation}
    \mathbf{X}_{t-L+1:t} \in \mathbb{R}^{L \times C}
\end{equation}
denote an input window of length $L$ with $C$ variables.
The goal is to predict
\begin{equation}
    \mathbf{Y}_{t+1:t+H} \in \mathbb{R}^{H \times C_o},
\end{equation}
where $H$ is the forecast horizon and $C_o$ is the output dimension.
For the GLD study, $C=C_o=1$ and the target is the daily \texttt{close} column.
The evaluated horizon set is
\begin{equation}
    H \in \{1, 3, 5, 7, 14, 30, 60, 90, 120, 180, 270, 360\}.
\end{equation}
These horizons cover short-term, middle-term, and long-term forecasting behavior.

This manuscript treats nonstationarity as a time-varying relationship between input and output windows.
Let $\mathcal{D}_t$ denote the local joint distribution around time $t$.
The nonstationary setting can be written as
\begin{equation}
    \mathcal{D}_t(\mathbf{X}_{t-L+1:t}, \mathbf{Y}_{t+1:t+H})
    \neq
    \mathcal{D}_{t+\tau}(\mathbf{X}_{t+\tau-L+1:t+\tau}, \mathbf{Y}_{t+\tau+1:t+\tau+H}).
\end{equation}
Operationally, the proposed models target window-level changes in location, scale, endpoint drift, range, and first-difference volatility.
These statistics do not exhaust all possible forms of nonstationarity, but they correspond to the shifts explicitly exposed to the adaptive normalization and gating modules.

\subsection{Mamba Backbone}

State space models provide an alternative to attention for sequence modeling.
The continuous-time form of a linear state space model is
\begin{equation}
    \mathbf{h}'(t)
    =
    \mathbf{A}\mathbf{h}(t)
    +
    \mathbf{B}\mathbf{x}(t),
    \qquad
    \mathbf{y}(t)
    =
    \mathbf{C}\mathbf{h}(t),
\end{equation}
where $\mathbf{h}(t)$ is the latent state, $\mathbf{x}(t)$ is the input, and $\mathbf{y}(t)$ is the output.
After discretization, the recurrence can be written abstractly as
\begin{equation}
    \mathbf{h}_t
    =
    \bar{\mathbf{A}}\mathbf{h}_{t-1}
    +
    \bar{\mathbf{B}}\mathbf{x}_t,
    \qquad
    \mathbf{y}_t
    =
    \mathbf{C}\mathbf{h}_t.
\end{equation}
Mamba introduces selective state spaces in which the discretization and readout parameters depend on the current input token \cite{mamba}.
The input-dependent update is summarized as
\begin{equation}
    \mathbf{h}_t
    =
    \bar{\mathbf{A}}_t\mathbf{h}_{t-1}
    +
    \bar{\mathbf{B}}_t\mathbf{x}_t,
    \qquad
    \mathbf{y}_t
    =
    \mathbf{C}_t\mathbf{h}_t.
\end{equation}
This selectivity is useful for forecasting because the model can adapt what information is retained across tokens.
Figure~\ref{fig:mamba-arch} summarizes the Mamba block used as the conceptual sequence-modeling backbone.

\begin{figure}[t]
    \centering
    \begin{minipage}[t]{0.46\linewidth}
\centering
\textbf{A. Block outline}\\[0.4em]
\resizebox{\linewidth}{!}{%
\begin{tikzpicture}[
    node distance=0.54cm and 0.42cm,
    >=Latex,
    arr/.style={->, draw=#1, line width=0.6pt},
    arr/.default={mambaLine},
    block/.style={
        draw=mambaLine,
        rounded corners=2pt,
        fill=white,
        align=center,
        font=\scriptsize,
        inner xsep=4pt,
        inner ysep=3pt,
        minimum height=0.58cm,
        text width=2.10cm
    },
    data/.style={block, fill=mambaBlue!18},
    ssm/.style={block, fill=mambaGreen!14, draw=mambaGreen!70!black},
    gate/.style={block, fill=mambaPurple!12, draw=mambaPurple!75!black},
    op/.style={block, fill=mambaGray!9}
]
\definecolor{mambaBlue}{HTML}{EAF4FF}
\definecolor{mambaGreen}{HTML}{DFF2E1}
\definecolor{mambaPurple}{HTML}{EFE8FF}
\definecolor{mambaGray}{HTML}{F4F6F8}
\definecolor{mambaLine}{HTML}{40556B}

\node[data] (input) {Input sequence\\$\mathbf{x}_t$};
\node[op, below=of input] (proj) {Input projection\\expand and split};
\node[ssm, below left=0.55cm and 0.18cm of proj] (ssmstream) {SSM stream\\conv $\rightarrow$ SiLU $\rightarrow$ scan};
\node[gate, below right=0.55cm and 0.18cm of proj] (gatestream) {Gate stream\\$\mathbf{z}_t=\operatorname{SiLU}(W_z\mathbf{x}_t)$};
\node[op, below=1.08cm of proj] (fusion) {Elementwise fusion\\$\mathbf{o}_t=\mathbf{y}_t\odot\mathbf{z}_t$};
\node[op, below=of fusion] (outproj) {Output projection};
\node[data, below=of outproj] (output) {Block output};

\draw[arr] (input) -- (proj);
\draw[arr] (proj) -- (ssmstream);
\draw[arr] (proj) -- (gatestream);
\draw[arr=mambaGreen!70!black] (ssmstream) |- (fusion);
\draw[arr=mambaPurple!75!black] (gatestream) |- (fusion);
\draw[arr] (fusion) -- (outproj);
\draw[arr] (outproj) -- (output);
\end{tikzpicture}%
}
\end{minipage}
\hfill
\begin{minipage}[t]{0.50\linewidth}
\centering
\textbf{B. Selective SSM component}\\[0.4em]
\resizebox{\linewidth}{!}{%
\begin{tikzpicture}[
    node distance=0.48cm and 0.30cm,
    >=Latex,
    arr/.style={->, draw=#1, line width=0.6pt},
    arr/.default={mambaLine},
    block/.style={
        draw=mambaLine,
        rounded corners=2pt,
        fill=white,
        align=center,
        font=\scriptsize,
        inner xsep=4pt,
        inner ysep=3pt,
        minimum height=0.58cm,
        text width=2.10cm
    },
    param/.style={block, fill=mambaGold!18, draw=mambaGold!80!black, text width=1.55cm},
    state/.style={block, fill=mambaGold!16, draw=mambaGold!80!black, text width=1.72cm},
    ssm/.style={block, fill=mambaGreen!14, draw=mambaGreen!70!black, text width=2.95cm},
    gate/.style={block, fill=mambaPurple!12, draw=mambaPurple!75!black, text width=1.45cm},
    data/.style={block, fill=mambaBlue!18, text width=2.35cm}
]
\definecolor{mambaBlue}{HTML}{EAF4FF}
\definecolor{mambaGreen}{HTML}{DFF2E1}
\definecolor{mambaPurple}{HTML}{EFE8FF}
\definecolor{mambaGold}{HTML}{FFF3C9}
\definecolor{mambaLine}{HTML}{40556B}

\node[data] (features) {SSM-stream features};
\node[block, below=of features] (heads) {Input-dependent\\parameter heads};
\node[param, below left=0.56cm and 0.30cm of heads] (delta) {Step size\\$\Delta_t$};
\node[param, below=0.56cm of heads] (bmap) {Input map\\$\mathbf{B}_t$};
\node[param, below right=0.56cm and 0.30cm of heads] (cout) {Readout\\$\mathbf{C}_t$};
\node[ssm, below=0.60cm of bmap] (disc) {Discretization\\$\bar{\mathbf{A}}_t=\exp(\Delta_t\mathbf{A})$\\$\bar{\mathbf{B}}_t=f(\Delta_t,\mathbf{A},\mathbf{B}_t)$};
\node[state, below left=0.55cm and 0.36cm of disc] (prev) {Previous state\\$\mathbf{h}_{t-1}$};
\node[ssm, below=0.55cm of disc] (update) {Selective state update\\$\mathbf{h}_t=\bar{\mathbf{A}}_t\mathbf{h}_{t-1}+\bar{\mathbf{B}}_t\mathbf{x}_t$};
\node[ssm, below=of update] (readout) {State readout\\$\mathbf{y}_t=\mathbf{C}_t\mathbf{h}_t$};
\node[gate, right=0.56cm of readout] (gate) {Gate\\$\mathbf{z}_t$};
\node[data, below=of readout] (modout) {Modulated output\\$\mathbf{o}_t=\mathbf{y}_t\odot\mathbf{z}_t$};

\draw[arr] (features) -- (heads);
\draw[arr] (heads) -- (delta);
\draw[arr] (heads) -- (bmap);
\draw[arr] (heads) -- (cout);
\draw[arr=mambaGold!80!black] (delta) -- (disc);
\draw[arr=mambaGold!80!black] (bmap) -- (disc);
\draw[arr=mambaGreen!70!black] (disc) -- (update);
\draw[arr=mambaGold!80!black] (prev) -- (update);
\draw[arr=mambaGreen!70!black] (update) -- (readout);
\draw[arr=mambaGold!80!black] (cout) |- (readout);
\draw[arr=mambaGreen!70!black] (readout) -- (modout);
\draw[arr=mambaPurple!75!black] (gate) |- (modout);
\end{tikzpicture}%
}
\end{minipage}

    \caption{Mamba block used as the sequence-modeling backbone. Panel A shows the high-level block structure with projection, selective state-space processing, gating, fusion, and output projection. Panel B expands the selective state-space component, where input-dependent parameters control the recurrent update and readout.}
    \label{fig:mamba-arch}
\end{figure}

\subsection{S-Mamba Forecasting Encoder}

Plain Mamba is a sequence block rather than a complete forecasting model.
Given a token sequence $\mathbf{U}\in\mathbb{R}^{B\times T\times d}$, a Mamba block produces another sequence
\begin{equation}
    \mathbf{V}=\mathcal{M}(\mathbf{U})\in\mathbb{R}^{B\times T\times d},
\end{equation}
with the same token axis $T$.
For generic sequence modeling, the token axis is usually time.
S-Mamba changes this setup for forecasting by using reversible window normalization and inverted variable tokens.

For an input window $\mathbf{X}\in\mathbb{R}^{B\times L\times C}$, S-Mamba first applies reversible normalization:
\begin{equation}
    \mu_{b,c}
    =
    \frac{1}{L}\sum_{t=1}^{L}X_{b,t,c},
    \qquad
    \sigma_{b,c}
    =
    \sqrt{
        \frac{1}{L}\sum_{t=1}^{L}(X_{b,t,c}-\mu_{b,c})^2 + \epsilon
    },
\end{equation}
\begin{equation}
    Z_{b,t,c}
    =
    \frac{X_{b,t,c}-\mu_{b,c}}{\sigma_{b,c}}.
\end{equation}
Each variable history is then embedded as one token:
\begin{equation}
    \mathbf{e}_{b,c}
    =
    W_{\mathrm{emb}}\mathbf{z}_{b,:,c}+ \mathbf{b}_{\mathrm{emb}},
    \qquad
    \mathbf{E}\in\mathbb{R}^{B\times C\times d_{\mathrm{model}}},
\end{equation}
where $\mathbf{z}_{b,:,c}\in\mathbb{R}^{L}$ is the normalized history of variable $c$.
Thus, S-Mamba scans along variable tokens whose features are full lookback histories, rather than scanning directly along time tokens.

Each encoder layer is bidirectional.
Let $\mathbf{E}^{(\ell-1)}$ be the input to layer $\ell$.
The forward and backward Mamba branches are
\begin{equation}
    \overrightarrow{\mathbf{H}}^{(\ell)}
    =
    \mathcal{M}_{f}^{(\ell)}(\mathbf{E}^{(\ell-1)}),
    \qquad
    \overleftarrow{\mathbf{H}}^{(\ell)}
    =
    \operatorname{rev}
    \left(
    \mathcal{M}_{b}^{(\ell)}
    \left(
    \operatorname{rev}(\mathbf{E}^{(\ell-1)})
    \right)
    \right),
\end{equation}
and the two directions are added:
\begin{equation}
    \mathbf{U}^{(\ell)}
    =
    \overrightarrow{\mathbf{H}}^{(\ell)}
    +
    \overleftarrow{\mathbf{H}}^{(\ell)}.
\end{equation}
The implemented layer then applies residual connections, layer normalization, and a feed-forward network:
\begin{equation}
    \tilde{\mathbf{E}}^{(\ell)}
    =
    \operatorname{LN}_{1}
    \left(
    \mathbf{E}^{(\ell-1)}+\mathbf{U}^{(\ell)}
    \right),
\end{equation}
\begin{equation}
    \mathbf{E}^{(\ell)}
    =
    \operatorname{LN}_{2}
    \left(
    \tilde{\mathbf{E}}^{(\ell)}+
    \operatorname{FFN}
    \left(
    \tilde{\mathbf{E}}^{(\ell)}
    \right)
    \right).
\end{equation}
After the final encoder layer, each variable token is projected to the horizon:
\begin{equation}
    \mathbf{P}_{b,c,:}
    =
    W_{\mathrm{out}}\mathbf{E}^{(L_e)}_{b,c,:}
    + \mathbf{b}_{\mathrm{out}},
    \qquad
    \mathbf{P}\in\mathbb{R}^{B\times C\times H}.
\end{equation}
The forecast is permuted back to the time-major layout and de-normalized:
\begin{equation}
    \hat{Y}_{b,h,c}
    =
    \sigma_{b,c} P_{b,c,h} + \mu_{b,c}.
\end{equation}

Figure~\ref{fig:s-mamba-encoder} summarizes the S-Mamba forecasting encoder.

\begin{figure}[t]
    \centering
    \centering
\resizebox{0.72\linewidth}{!}{%
\begin{tikzpicture}[
    node distance=0.66cm and 0.42cm,
    >=Latex,
    arr/.style={->, draw=#1, line width=0.6pt},
    arr/.default={mLine},
    block/.style={
        draw=mLine,
        rounded corners=2pt,
        fill=white,
        align=center,
        font=\small,
        inner xsep=5pt,
        inner ysep=4pt,
        minimum height=0.72cm,
        text width=2.70cm
    },
    data/.style={block, fill=mBlue!16},
    norm/.style={block, fill=mGold!16, draw=mGold!80!black},
    enc/.style={block, fill=mGreen!14, draw=mGreen!70!black, text width=3.60cm},
    proj/.style={block, fill=mPurple!12, draw=mPurple!75!black}
]
\definecolor{mBlue}{HTML}{EAF4FF}
\definecolor{mGreen}{HTML}{DFF2E1}
\definecolor{mPurple}{HTML}{EFE8FF}
\definecolor{mGold}{HTML}{FFF3C9}
\definecolor{mLine}{HTML}{40556B}

\node[data] (input) {Input window\\$\mathbf{X}\in\mathbb{R}^{B\times L\times C}$};
\node[norm, below=of input] (norm) {Window norm\\time-axis mean/std};
\node[proj, below=of norm] (embed) {Inverted embedding\\$\mathbf{E}\in\mathbb{R}^{B\times C\times d}$};
\node[enc, below left=0.62cm and 0.20cm of embed] (forward) {Forward Mamba\\$\overrightarrow{\mathbf{H}}=\mathcal{M}_f(\mathbf{E})$};
\node[enc, below right=0.62cm and 0.20cm of embed] (backward) {Backward Mamba\\$\overleftarrow{\mathbf{H}}=\operatorname{rev}(\mathcal{M}_b(\operatorname{rev}(\mathbf{E})))$};
\node[enc, below=1.18cm of embed] (fusion) {Bidirectional encoder layer\\residual + layer norm + FFN};
\node[proj, below=of fusion] (head) {Linear horizon head\\$\mathbf{P}\in\mathbb{R}^{B\times C\times H}$};
\node[norm, below=of head] (denorm) {De-normalization};
\node[data, below=of denorm] (forecast) {Forecast\\$\hat{\mathbf{Y}}\in\mathbb{R}^{B\times H\times C}$};

\draw[arr] (input) -- (norm);
\draw[arr] (norm) -- (embed);
\draw[arr] (embed) -- (forward);
\draw[arr] (embed) -- (backward);
\draw[arr=mGreen!70!black] (forward) |- (fusion);
\draw[arr=mGreen!70!black] (backward) |- (fusion);
\draw[arr] (fusion) -- (head);
\draw[arr] (head) -- (denorm);
\draw[arr] (denorm) -- (forecast);
\end{tikzpicture}%
}

    \caption{S-Mamba forecasting encoder: reversible normalization, inverted variable-token embedding, bidirectional Mamba encoding, horizon projection, and de-normalization back to the original price scale.}
    \label{fig:s-mamba-encoder}
\end{figure}

\subsection{Nonstationarity Statistics}

The NS-Mamba variants expose window statistics to two different modules.
Adaptive normalization uses raw window statistics to generate affine parameters, while shift-aware gates use standardized statistics to produce fusion weights.
Both are computed from the raw input window before normalization.

For adaptive normalization, the implementation constructs a ten-dimensional statistic vector for each sample and variable:
\begin{equation}
    \psi_{b,c}
    =
    [
    \mu_{b,c},
    \sigma_{b,c},
    X_{b,1,c},
    X_{b,L,c},
    \mu^{r}_{b,c},
    \sigma^{r}_{b,c},
    X_{b,L,c}-X_{b,1,c},
    \min_t X_{b,t,c},
    \max_t X_{b,t,c},
    \nu_{b,c}
    ],
\end{equation}
where $\mu_{b,c}$ and $\sigma_{b,c}$ are full-window mean and standard deviation, $\mu^{r}_{b,c}$ and $\sigma^{r}_{b,c}$ are computed on the most recent segment of the window, and $\nu_{b,c}$ is the standard deviation of first differences.
The recent segment length is controlled by the hyperparameter \texttt{ns\_recent\_ratio}, with default value 0.25.

For shift-aware fusion, the implementation uses an eight-dimensional standardized feature vector:
\begin{equation}
    \phi_{b,c}
    =
    [
    \delta_{\mu},
    \delta_{\sigma},
    s,
    \bar{\nu},
    \delta_{\nu},
    r,
    z_L,
    z_1
    ]_{b,c}.
\end{equation}
Let $\sigma_{b,c}^{+}=\sigma_{b,c}+\epsilon$.
The entries are
\begin{align}
    \delta_{\mu}
    &=
    \frac{\mu^{r}_{b,c}-\mu_{b,c}}{\sigma_{b,c}^{+}},
    &
    \delta_{\sigma}
    &=
    \log
    \frac{\sigma^{r}_{b,c}+\epsilon}{\sigma_{b,c}^{+}},
    \\
    s
    &=
    \frac{X_{b,L,c}-X_{b,1,c}}{\sigma_{b,c}^{+}},
    &
    \bar{\nu}
    &=
    \frac{\nu_{b,c}}{\sigma_{b,c}^{+}},
    \\
    \delta_{\nu}
    &=
    \log
    \frac{\nu^{r}_{b,c}+\epsilon}{\nu_{b,c}+\epsilon},
    &
    r
    &=
    \frac{\max_t X_{b,t,c}-\min_t X_{b,t,c}}{\sigma_{b,c}^{+}},
    \\
    z_L
    &=
    \frac{X_{b,L,c}-\mu_{b,c}}{\sigma_{b,c}^{+}},
    &
    z_1
    &=
    \frac{X_{b,1,c}-\mu_{b,c}}{\sigma_{b,c}^{+}}.
\end{align}
Here $\nu^{r}_{b,c}$ is the recent-segment standard deviation of first differences.
The fusion gates project $\phi_{b,c}$ into the hidden dimension and combine it with the hidden representation when one is available.
When the number of statistic tokens does not match the number of target tokens, the statistic context is averaged over variables and broadcast to the required token count.

\subsection{Adaptive Normalization}

Adaptive normalization replaces fixed reversible normalization with a learned affine transformation conditioned on $\psi_{b,c}$.
For each variable token, an MLP maps $\psi_{b,c}$ to raw affine parameters, and the bounded transformation is
\begin{equation}
    \gamma_{b,c}
    =
    1 + a \tanh(g_{\gamma}(\psi_{b,c})),
    \qquad
    \beta_{b,c}
    =
    a \tanh(g_{\beta}(\psi_{b,c})),
\end{equation}
where $a$ is the adaptive-scale hyperparameter.
The normalized input is
\begin{equation}
    \tilde{X}_{b,t,c}
    =
    \gamma_{b,c}
    \frac{X_{b,t,c}-\mu_{b,c}}{\sigma_{b,c}+\epsilon}
    +
    \beta_{b,c}.
\end{equation}
The same cached $\gamma_{b,c}$, $\beta_{b,c}$, $\mu_{b,c}$, and $\sigma_{b,c}$ are used to invert the transformation at the output:
\begin{equation}
    \hat{Y}_{b,h,c}
    =
    \left(
    \frac{\tilde{Y}_{b,h,c}-\beta_{b,c}}{\gamma_{b,c}}
    \right)
    (\sigma_{b,c}+\epsilon)
    +
    \mu_{b,c}.
\end{equation}

\subsection{Series Decomposition}

The decomposition module separates the normalized window into low-frequency, medium-frequency, and residual components using moving averages.
Let $\operatorname{MA}_{m}(\cdot)$ denote a centered moving average with boundary replication and window size $m$.
The implementation computes
\begin{align}
    \mathbf{Z}^{\mathrm{tr}}
    &=
    \operatorname{MA}_{m}(\mathbf{Z}),
    \\
    \mathbf{R}^{0}
    &=
    \mathbf{Z}-\mathbf{Z}^{\mathrm{tr}},
    \\
    \mathbf{Z}^{\mathrm{mf}}
    &=
    \operatorname{MA}_{m}(\mathbf{R}^{0}),
    \\
    \mathbf{Z}^{\mathrm{re}}
    &=
    \mathbf{R}^{0}-\mathbf{Z}^{\mathrm{mf}}.
\end{align}
The superscripts denote low-frequency trend-like, medium-frequency, and residual components.
This decomposition is an architectural bias rather than a claim that GLD prices are generated by a perfect additive seasonal process.

\subsection{Shift-Aware Fusion}

Shift-aware gates use $\phi_{b,c}$ to control either branch fusion or hidden-state interpolation.
For decomposition branches with horizon-level outputs $\mathbf{U}_q\in\mathbb{R}^{B\times C\times H}$, where $q\in\{\mathrm{tr},\mathrm{mf},\mathrm{re}\}$, the branch gate computes
\begin{equation}
    \boldsymbol{\alpha}_{b,c,:,h}
    =
    \operatorname{softmax}_{q}
    \left(
    g_{\alpha}(\phi_{b,c})
    \right),
\end{equation}
and fuses the branch outputs as
\begin{equation}
    \mathbf{U}_{b,c,h}
    =
    \sum_q \alpha_{b,c,q,h}\mathbf{U}_{q,b,c,h}.
\end{equation}
The same gate class is also used for scale fusion in the multi-scale models, where the softmax dimension ranges over temporal scales rather than decomposition components.
When no decomposition branch is present, a sigmoid gate can interpolate between the encoded hidden state and the input embedding before horizon projection:
\begin{equation}
    \tilde{\mathbf{E}}
    =
    \mathbf{G}\odot \mathbf{E}_{\mathrm{enc}}
    +
    (1-\mathbf{G})\odot \mathbf{E}_{\mathrm{in}}.
\end{equation}

\subsection{Variant Definitions}

The Mamba-family variant names encode which modules are active.
\textbf{A} denotes adaptive normalization, \textbf{D} denotes decomposition, and \textbf{G} denotes shift-aware gated fusion.
\textbf{M} denotes multi-scale temporal branches, \textbf{S} denotes an NLinear-style residual skip, and \textbf{P} denotes patch-temporal Mamba branches.
All variants produce forecasts with shape $\mathbb{R}^{B \times H \times C}$ and use reversible output scaling when normalization is enabled.

Let $\mathcal{N}(\cdot)$ denote standard reversible normalization and $\mathcal{F}_{\theta}(\cdot)$ denote one inverted S-Mamba forecasting branch.
The S-Mamba baseline is
\begin{equation}
    \hat{\mathbf{Y}}_{\mathrm{S}}
    =
    \mathcal{N}^{-1}
    \left(
    \mathcal{F}_{\theta}(\mathcal{N}(\mathbf{X}))
    \right).
\end{equation}
NS-Mamba-A replaces $\mathcal{N}$ with adaptive normalization.
NS-Mamba-D applies standard reversible normalization, decomposes the normalized signal, and sums three independently encoded component forecasts:
\begin{equation}
    \hat{\mathbf{Y}}_{\mathrm{D}}
    =
    \mathcal{N}^{-1}
    \left(
    \sum_{q \in \{\mathrm{tr},\mathrm{mf},\mathrm{re}\}}
    \mathcal{F}_{\theta_q}(\mathbf{Z}^{q})
    \right).
\end{equation}
NS-Mamba-AD combines adaptive normalization with decomposition, and NS-Mamba-ADG adds the shift-aware branch gate.

The MS-Mamba family adds temporal scale branches.
Let $\mathcal{P}_{s}(\cdot)$ be average-pooling downsampling by scale $s$, and let $\mathcal{S}$ be a set of scales such as $\{1,2,4\}$.
For a decomposition-enabled multi-scale model, each scale produces
\begin{equation}
    \mathbf{U}_{s}
    =
    \sum_{q \in \{\mathrm{tr},\mathrm{mf},\mathrm{re}\}}
    \mathcal{F}_{\theta_{s,q}}
    \left(
    \mathcal{P}_{s}(\mathbf{Z}^{q})
    \right).
\end{equation}
MS-Mamba-MD uses mean scale fusion:
\begin{equation}
    \hat{\mathbf{Y}}_{\mathrm{MD}}
    =
    \mathcal{N}^{-1}
    \left(
    \frac{1}{|\mathcal{S}|}
    \sum_{s \in \mathcal{S}}
    \mathbf{U}_{s}
    \right),
\end{equation}
while MS-Mamba-MDG replaces the mean with a shift-aware softmax over scales.
MS-Mamba-M and MS-Mamba-MG are the corresponding multi-scale variants without decomposition.

MS-Mamba-MDS adds a last-value residual linear skip:
\begin{equation}
    \mathbf{U}_{\mathrm{skip}}
    =
    \mathbf{x}_{L}
    +
    \mathbf{W}_{\mathrm{skip}}
    (\mathbf{Z}-\mathbf{x}_{L}),
    \qquad
    \hat{\mathbf{Y}}_{\mathrm{MDS}}
    =
    \mathcal{N}^{-1}
    \left(
    \lambda \mathbf{U}_{\mathrm{MD}}
    +
    (1-\lambda)\mathbf{U}_{\mathrm{skip}}
    \right),
\end{equation}
where $\lambda=\sigma(\rho)$ is learned.
PMS-Mamba replaces inverted variable-token branches with patch-temporal Mamba branches.
For patch length $p$ and stride $r$, a scale branch forms patch tokens
\begin{equation}
    \mathbf{e}_{j}^{(s)}
    =
    \mathbf{W}_{p}
    \operatorname{vec}
    \left(
    \mathcal{P}_{s}(\mathbf{Z})_{j:j+p-1}
    \right),
\end{equation}
then applies Mamba over the patch index $j$.
PMS-Mamba-PD is the patch-temporal counterpart of MS-Mamba-MD.

\begin{table}[t]
\centering
\caption{Meaning of the Mamba-family variant names used in this manuscript.}
\label{tab:mamba-variant-meaning}
{\small
\begin{tabular}{p{0.19\linewidth}p{0.27\linewidth}p{0.16\linewidth}p{0.26\linewidth}}
\toprule
Variant & Active modules & Fusion & Purpose \\
\midrule
NS-Mamba-A & Adaptive norm & Single path & Test learned normalization \\
NS-Mamba-D & Decomposition & Sum & Test component separation \\
NS-Mamba-AD & Adaptive norm + decomposition & Sum & Combine A and D \\
NS-Mamba-ADG & Adaptive norm + decomposition & Shift gate & Test gated component fusion \\
MS-Mamba-M & Multi-scale & Mean & Test temporal scales alone \\
MS-Mamba-MG & Multi-scale & Shift gate & Test gated scale fusion \\
MS-Mamba-MD & Multi-scale + decomposition & Mean & Promoted Mamba model \\
MS-Mamba-MDG & Multi-scale + decomposition & Shift gate & Test gated multi-scale decomposition \\
MS-Mamba-MDS & Multi-scale + decomposition + skip & Mean + skip & Test NLinear-style residual path \\
PMS-Mamba-PD & Patch Mamba + decomposition & Mean & Test patch Mamba \\
\bottomrule
\end{tabular}
}
\end{table}

\subsection{Forecast Head and Loss}

The forecast head projects the final hidden representation to $H$ future values.
The training objective is mean squared error over the forecast horizon:
\begin{equation}
    \mathcal{L}
    =
    \frac{1}{BHC_o}
    \sum_{b=1}^{B}
    \sum_{h=1}^{H}
    \sum_{c=1}^{C_o}
    \left(
    \hat{Y}_{b,h,c} - Y_{b,h,c}
    \right)^2.
\end{equation}
Inverse transformation is enabled by default in the reported GLD runs, so saved predictions, targets, and metrics are reported on the original GLD price scale unless a run uses \texttt{--no-inverse}.
